\documentclass[10pt,a4paper]{amsart}
\usepackage[margin=19mm]{geometry}
\usepackage{amsmath,amssymb,mathtools}
\usepackage[scr=rsfs,cal=euler]{mathalfa}
\usepackage{xfrac}
\usepackage[unicode,hidelinks,bookmarks=false]{hyperref}
\newcommand{\ie}{i.e.\ }
\newcommand{\cf}{cf.\ }
\newcommand{\resp}{respectively\ }
\newcommand{\Subsection}{\subsection}
\providecommand{\nobreakdash}{\nobreak\hbox{-}}
\setlength{\emergencystretch}{2em}
\multlinegap0pt
\DeclareMathOperator{\Cov}{\mathrm{Cov}}
\DeclareMathOperator{\Gal}{\mathrm{Gal}}
\DeclareMathOperator{\im}{\mathrm{im}}
\DeclareMathOperator{\loc}{\mathrm{loc}}
\theoremstyle{plain}
\newtheorem{theorem}{Theorem}
\newtheorem{proposition}[theorem]{Proposition}
\newtheorem{lemma}[theorem]{Lemma}
\newtheorem{corollary}[theorem]{Corollary}
\theoremstyle{definition}
\newtheorem{definition}[theorem]{Definition}
\newtheorem{assumption}[theorem]{Assumption}
\theoremstyle{remark}
\newtheorem{remark}[theorem]{Remark}
\title[Rigidity and Galois theory of 3-manifolds branched over Chebotarev links]{Rigidity and Galois theory of 3-manifolds branched over Chebotarev links — Theorem 6.7}
\author{Nadav Gropper, Jun Ueki, Yi Wang}
\date{Source: arXiv:2604.09469v2 (CC BY 4.0)}
\begin{document}
\maketitle

\noindent\emph{Source and licence.} The delivered work is the article shown in the title above, version arXiv:2604.09469v2 (CC BY 4.0), distributed under the Creative Commons Attribution 4.0 International licence, as recorded in the arXiv licence metadata for this exact version and shown on the abstract page saved with this item. The full licence notice, which carries the licence URL and the address of that abstract page, is the bundled file \texttt{licenses/arxiv-2604.09469-CC-BY-4.0.txt}.\par\smallskip

\noindent\textit{Editorial scope.} Counted unit: the article's Theorem 6.7 (source0.tex lines 1008--1014, source label \texttt{thm:exact}); statement and proof are the article's own text line for line, with no line omitted, substituted or re-flowed. Required context retained uncounted: the article's setup blocks the counted statement and proof invoke (source0.tex lines 509--511, 965--970 and 986--992). Editorial formatting only: an AMS \texttt{amsart} preamble that restates the article's own operator definitions (\texttt{\textbackslash Cov}, \texttt{\textbackslash Gal}, \texttt{\textbackslash im}, \texttt{\textbackslash loc}) and the theorem environments the retained lines use, including the article's own \texttt{mainthm} environment name for its main theorem, which the retained context declares but does not use; the article ships no class or style file; sequential numbering of the retained environments in order of appearance, whereas the article prints the counted theorem as Theorem 6.7; equation numbering is not reproduced and every cross-reference inside the retained text resolves to the wrapper's numbering. No citation occurs inside any retained line, so the wrapper carries no bibliography. No asset-inclusion command, no drawing environment and no image asset occurs in this package.


\section*{Retained context: cor:incompressible2 (source0.tex lines 509--511)}

\begin{corollary}\label{cor:incompressible2}
Let $\mathcal{L} \subset S^3$ be Chebotarev, let $h: M \to S^3$ be a cover in $\Cov(S^3, \mathcal L)$, and let $L$ be a finite sublink of $\mathcal L_M$. Then there is a finite sublink $L'$ containing $L$ such that the exterior of $L'$ in $M$ is aspherical and all its boundary tori are $\pi_1$-injective.
\end{corollary}


\section*{Retained context: lma:localduality (source0.tex lines 965--970)}

\begin{lemma}\label{lma:localduality}
We have the following isomorphism, coming from the local duality for decomposition groups:
\begin{equation}
    \Xi^1: P^1(\Gal(M, \mathcal{L}_M), A) \to P^1(\Gal(M, \mathcal{L}_M), A)^\vee.
\end{equation}
\end{lemma}


\section*{Retained context: lma:lefschetz (source0.tex lines 986--992)}

\begin{lemma}\label{lma:lefschetz}
Let $X = M \setminus \bigcup_{K \in L}V_K^{\circ}$ be a compact oriented aspherical 3-manifold with boundary $\partial X = \bigcup_{K \in L}\partial V_K$ a finite disjoint union of $\pi_1$-injective tori. Then the sequence
\begin{equation}\label{eq:58}
    H^1(\Gal(M, L); A) \to \prod_{K \in L}H^1(D_{\overline{K}|K}; A) \to H^1(\Gal(M, L); A)^\vee
\end{equation}
is exact. The first map is given by $\loc_L$, the second map is given by $(\loc_L)^\vee \circ \Xi^1$, and $\Xi^1$ is the Poincar\'e duality pairing on $\partial X$.
\end{lemma}


\section{The counted unit: Theorem 6.7 and its complete original proof}

\begin{theorem}\label{thm:exact}
Let $h: M \to S^3 \in \Cov(S^3, \mathcal{L})$. The sequence
\begin{equation}
    H^1(\Gal(M, \mathcal{L}_M), A) \to P^1(\Gal(M, \mathcal{L}_M), A) \to H^1(\Gal(M, \mathcal{L}_M), A)^\vee
\end{equation}
is exact. Denote the first map by $\loc_{\mathcal{L}}$; the second map is given by $(\loc_\mathcal{L})^\vee \circ \Xi^1$.
\end{theorem}

\begin{proof}
By Corollary \ref{cor:incompressible2}, for any finite sublink $L$ we can add components so that all boundary components are $\pi_1$-injective and aspherical; Lemma \ref{lma:lefschetz} then gives the desired exact sequence. We now pass to the full restricted product $P^1 = P^1(\Gal(M, \mathcal{L}_M), A)$ over all of $\mathcal{L}_M$. The exactness of this sequence is equivalent to showing that $\im(\loc_{\mathcal{L}})$ equals its own orthogonal complement in $P^1$ under the pairing
\begin{equation}
    P^1 \times P^1 \rightarrow \bigoplus H^2(D_{\overline{K}|K}, A) \to \mathbb{Q}/\mathbb{Z}
\end{equation}
where the first map is the componentwise cup product and the second map is summation.

\medskip

\underline{Claim 1: $\im(\loc_{\mathcal{L}}) \subset \im(\loc_{\mathcal{L}})^\perp$.}

\medskip

Given two classes $x, y \in H^1(\Gal(M, \mathcal{L}_M), A)$, the local cup product $(\loc_{\mathcal{L}}(x))_K \cup (\loc_{\mathcal{L}}(y))_K \in H^2(D_{\overline{K}|K}, A)$ can be nontrivial only if at least one of $x_{D_{\overline{K}|K}}$ or $y_{D_{\overline{K}|K}}$ lies in the ramified part of $H^1(D_{\overline{K}|K}, A)$ (follows from self-orthogonality of $H^1_{ur}$, i.e., Lemma \ref{lma:localduality}). Since $x$ and $y$ are globally unramified outside a finite set of knots, this can only occur for finitely many $K$. Let $L \subset \mathcal{L}_M$ be a finite subset including all such $K$, chosen to have aspherical complement and $\pi_1$-injective boundary by Corollary \ref{cor:incompressible2}. Because $x$ and $y$ are unramified outside $L$, there exist classes $\overline{x}, \overline{y} \in H^1(\Gal(M, L), A)$ such that $x = \inf(\overline{x})$ and $y = \text{inf}(\overline{y})$. Thus, the global cup product evaluates as:
\begin{equation}
    \loc_\mathcal{L}(x) \cup \loc_\mathcal{L}(y) = \pi_L(\loc_{\mathcal{L}}(x)) \cup \pi_L(\loc_{\mathcal{L}}(y)) = \loc_L(\overline{x}) \cup \loc_L(\overline{y}) = 0
\end{equation}
where the last equality uses Lemma \ref{lma:lefschetz}, giving $\im(\loc_{\mathcal{L}}) \subset \im(\loc_{\mathcal{L}})^\perp$.

\medskip

\underline{Claim 2: $\im(\loc_{\mathcal{L}})$ is closed in $P^1$.}

\medskip

We now show that $\loc_{\mathcal{L}}$ is proper, i.e., the preimage of any compact subset of $P^1$ is compact in $H^1(\Gal(M, \mathcal{L}_M), A)$. A compact neighbourhood of the identity in $P^1$ is contained in 
\begin{equation}
    P_L = \prod_{K \in L}H^1(D_{\overline{K}|K}, A) \times \prod_{K \in \mathcal{L}_M \setminus L}H^1_{ur}(D_{\overline{K}|K}, A)
\end{equation}
for some finite $L$. Since $H^1(\Gal(M, \mathcal{L}_M), A)$ is discrete, it suffices to show that $\loc_{\mathcal{L}}^{-1}(P_L)$ is finite. Any $x \in \loc_{\mathcal{L}}^{-1}(P_L)$ satisfies $x|_{I_{\overline{K}|K}} = 0$ for all $K \notin L$, hence factors through the maximal quotient of $\Gal(M, \mathcal{L}_M)$ unramified outside $L$, which is $\Gal(M, L) = \widehat{\pi}_1(M \setminus L)$. Via the inflation-restriction exact sequence:
\begin{equation}
    0 \to H^1(\Gal(M, L), A) \to_{\text{inf}} H^1(\Gal(M, \mathcal{L}_M), A)
\end{equation}
we have $\loc_{\mathcal{L}}^{-1}(P_L) \subset \im(\text{inf})$ (particularly using that inflation is injective in $H^1$ for normal subgroups). Since the exterior of $L$ is a compact 3-manifold with boundary, $\pi_1(M \setminus L)$ is finitely presented, hence $\Gal(M, L)$ is topologically finitely generated, making $H^1(\Gal(M, L), A)$ finite. Therefore, $\loc_{\mathcal{L}}^{-1}(P_L)$ is finite, and $\loc_{\mathcal{L}}$ is proper. A proper map from a discrete space into a locally compact space has closed image, so $\im(\loc_{\mathcal{L}})$ is closed in $P^1$. 

\medskip

\underline{Claim 3: $\im(\loc_{\mathcal{L}})^\perp \subset \im(\loc_{\mathcal{L}})$.}

\medskip

Let $x \in P^1$ satisfy $x \cup \loc_{\mathcal{L}}(y) = 0$ for all $y \in H^1(\Gal(M, \mathcal{L}_M), A)$. For any finite $L$ containing the finite set of knots where $x$ is ramified and whose exterior is aspherical with $\pi_1$-injective boundary, projecting to $P^1(\mathcal L)$ gives $|pi_L(x) \in \im(\loc_L)^\perp = \im(\loc_L)$ by Lemma \ref{lma:lefschetz}. Here we use that every class in $H^1(\Gal(M, \mathcal L), A)$ is in the image of the inflation $H^1(\Gal(M, \mathcal L_M), A)$, so orthogonality against all global classes at the infinite level implies orthogonality against all classes at the finite level.
% Then $x - \loc_{\mathcal L}(y)$ vanishes on $L$ and is unramified outside $L$, i.e.
% \begin{equation}
%     x - \loc_{\mathcal L}(y) \in N_L := \{z \in P^1 \mid z_K = 0 \text{ for } K \subset L, z_K \in H^1_{\mathrm{ur}}(D_{\bar K \mid K}; A) \text{ for } K \not\subset L\}
% \end{equation}
% These $N_L$ form a neighborhood basis at zero in $P^1$. Thus $x \in \bigcap_L(\im(\loc_{\mathcal L}) + N_L) = \overline{\im\loc_{\mathcal L}} = \im(\loc_{\mathcal L})$ by Claim 2. 
Since $\im(\loc_{\mathcal{L}})$ is closed by Claim 2:
\begin{equation}
    x \in \bigcap_L\pi_L^{-1}(\im(\loc_L)) \subset \bigcap_L\pi_L^{-1}(\pi_L(\im(\loc_{\mathcal{L}}))) = \overline{\im(\loc_{\mathcal{L}})} = \im(\loc_{\mathcal{L}}).
\end{equation}
The first inclusion uses that $\im(\loc_L)$ at the finite level equals $\pi_L(\im(\loc_{\mathcal{L}}))$, i.e., every finite-level element in $\im(\loc_L)$ lifts to a global element in $\im(\loc_{\mathcal{L}})$. The third equality uses Claim 2.

\medskip

Claim 1 and 3 together give $\im(\loc_{\mathcal{L}}) = \im(\loc_{\mathcal{L}})^\perp$, which under the identification via $\Xi^1$ is equivalent to exactness of the stated sequence.
\end{proof}

\end{document}
